\subsection{Refitted graph readout}
The mean-aggregation comparator uses the original query-masked endpoint inputs, message encoder, donor-balanced fitting loss and learning-rate and early-stop rules. It uses mean pooling within each neighborhood and the original mean/max pooling across the endpoints. It omits explicit degree inputs and degree-scaled pooling. Its original GlaS four folds yield normal-link AUROC between 95.7\% and 98.2\% with better loss than a prevalence intercept in all four. All 280 query-bit flip checks give exact equality. The mean model can still use neighborhood content that covaries with degree or recognize an empty neighborhood.

\begin{table*}[t]
\centering
\caption{\textbf{Response after refitting with mean aggregation.} AUC and top-four recovery use absolute matched probability change and donor-macro aggregation. Top-four change is refit minus original in percentage points. Intervals are donor-paired. The tiers share donors.}
\label{tab:refit}
\small
\begin{tabular}{lrrrr}
\toprule
Switch tier & Donors & Response AUROC & Top four & Top-four change \\
\midrule
Original & 12 & 91.5 [89.1, 93.7] & 46.6 [40.5, 52.4] & -40.9 [-49.8, -32.8] \\
Extra length-matched & 11 & 92.7 [90.1, 95.1] & 45.3 [40.2, 50.6] & -37.6 [-48.1, -23.0] \\
Joint-matched & 11 & 92.9 [90.3, 95.3] & 51.4 [41.7, 59.7] & -35.2 [-45.3, -25.6] \\
\bottomrule
\end{tabular}

\end{table*}

The original and refitted networks use a query-excluded adjacency for all tokens. A changed observed bit can affect a loss through the scoring formula without moving the fitted probability. The pair MLP is an invariance control because geometry and roles stay fixed. The degree-only probe concentrates its changed features on the four target pairs by construction. It tests whether the fitted model responds to that input. Its perfect recovery and the fixed-weight clamp should be interpreted with the refit in \cref{tab:refit}. SpotTarget studies degree-sensitive bias from including and excluding target edges \cite{zhu2024spot}. Our paired construction exposes a residual input difference after exclusion.

\clearpage
